\reporttitle{总览}
\renewcommand{\stagefigureprefix}{0}\setcounter{figure}{0}

\section*{我们做了什么}
本项目在香山昆明湖V2处理器上新增了一条自定义向量点积指令vdot.vv（下文简称vdot）。它一次读取两个向量寄存器中的8个有符号8位整数，两两相乘后求和，把结果写入目的寄存器的最低32位，其余96位清零。点积是量化大模型推理中矩阵向量乘的核心运算；用标准RISC-V向量指令（RVV）完成同样的计算，需要宽乘、归约和多次配置切换，专用指令把这些步骤合并成了一条。

围绕这条指令，我们完成了：指令编码与语义定义；在香山已有的向量执行单元中新增点积功能单元（VDot）并完成RTL实现；在参考模型NEMU中实现同一语义，与RTL进行逐条指令的协同仿真；按正数、负数、边界值和随机输入分层验证；与等价的RVV程序比较周期，并评估面积和时序；提供可直接调用的C函数库。

\section*{最终交付的设计}
硬件上，VDot功能单元接入香山的向量执行单元VFEX4（文中称“父级”），与原有的整数除法、浮点除法/开方单元共享派发和写回，处理器原有的译码、重命名、调度和提交机制全部复用。功能单元提供两种实现：默认的单请求实现，以及可选的两槽流水实现。两槽实现用乘积槽P和结果槽R组成两级流水，接收请求后最早2拍给出结果，每拍可以接收一笔新请求。本报告的性能、面积和时序数据都来自两槽实现的同一个固定版本，文中称为“受测版本”。

软件上，提供普通C接口的8×32整数矩阵向量乘（GEMV）函数库。其中C8版本在8行计算之间复用向量配置，是本报告选用的调用实现。

扩展方面，十六项INT8点积指令和INT16元素的参数化模块都已完成功能验证，尚未纳入性能评估。

\section*{主要结果}
主要结果汇总于表\ref{tab:0-results}，每一项都注明了正文中的详细位置。

\begin{table}[!htbp]\centering
\small
\caption{主要结果汇总}\label{tab:0-results}
\begin{tabular}{>{\raggedright\arraybackslash}p{0.15\linewidth}>{\raggedright\arraybackslash}p{0.62\linewidth}>{\raggedright\arraybackslash}p{0.15\linewidth}}
\toprule
方面 & 结果 & 详见\\
\midrule
三类基础测试 & 正数（结果204）、负数（$-8$）和边界值（131072）的完整128位结果均与期望一致 & 第\ref{sec:2-verify}节\\
功能正确性 & 算术RTL的10022组输入（22组定向、10000组随机）全部与独立数学模型一致；全核117个用例、124次运行在协同仿真中与参考模型一致 & 第\ref{sec:2-verify}、\ref{sec:3-diff}节\\
错误检出 & 在参考模型的结果中人为翻转1位后，协同仿真能检出差异并终止 & 第\ref{sec:3-diff}节\\
流水性能 & 两槽实现接收后最早2拍给出结果，每拍可接收一笔新请求 & 第\ref{sec:2-datapath}节\\
微基准加速 & 相对等价RVV程序：独立批量3.07倍，依赖链2.93倍，单次点积1.20倍，含加载和存储的流式计算1.005倍 & 表\ref{tab:3-micro-results}\\
函数调用 & 8×32 GEMV：C8版本273周期，比较快的RVV实现（299周期）快1.095倍 & 第\ref{sec:3-c8}节\\
完整模型 & 在一个六层Transformer的完整前向计算中，VDot内核通过全部状态比较 & 第\ref{sec:3-forward}节\\
面积 & VDot功能单元17,298\,$\mu\mathrm{m}^{2}$，约占所在执行单元VFEX4的6.1\%（Nangate45工艺库） & 第\ref{sec:3-cost}节\\
时序 & 2\,ns约束下VDot功能单元全部路径满足要求；1\,ns约束下最差裕量为$-0.27$\,ns，仍需优化 & 第\ref{sec:3-cost}节\\
扩展能力 & 十六项点积通过6个全核差分用例；INT16参数化模块通过94组定向和10000组随机输入 & 第\ref{sec:2-dot16}、\ref{sec:2-int16}节\\
\bottomrule
\end{tabular}
\end{table}

从结果看，vdot的收益主要来自可以连续发起的点积：在独立批量和依赖链场景中，受测版本的周期都降到RVV程序的约三分之一。当程序的瓶颈在数据加载和存储时，单条指令的加速对整体帮助有限；在函数调用和完整模型层面，配置切换、数据搬运和循环开销也会摊薄收益，因此我们同时做了C8这类软件层面的优化。

\section*{报告结构与赛题要求的对应}
\begin{itemize}
\item 第一阶段（环境部署与向量执行过程）：环境搭建、操作日志、向量加法vadd的执行结果与波形，对应“阶段一交付物”。
\item 第二阶段（指令设计与基础测试）：指令语义、数据通路、两槽流水、RTL实现，以及正数、负数、边界值三类测试，对应“阶段二交付物”和“技术路径合理性”。十六项点积与INT16扩展见第\ref{sec:2-dot16}、\ref{sec:2-int16}节，对应“扩展应用潜力”。
\item 第三阶段（协同仿真与性能评估）：NEMU扩展、RTL与NEMU协同仿真、等价RVV基线、周期对比、C函数库、面积与时序，对应“阶段三交付物”“功能验证有效性”“软硬件协同理解”和“微架构优化分析”。
\end{itemize}

\section*{常用术语}
全文反复出现的术语见表\ref{tab:0-terms}。

\begin{table}[!htbp]\centering
\small
\caption{常用术语}\label{tab:0-terms}
\begin{tabular}{>{\raggedright\arraybackslash}p{0.17\linewidth}>{\raggedright\arraybackslash}p{0.75\linewidth}}
\toprule
术语 & 含义\\
\midrule
vdot、VDot & 本项目新增的点积指令；实现该指令的功能单元\\
RVV & RISC-V标准向量扩展；本报告的对比基线是用RVV指令写成的等价程序\\
VLEN、SEW、VL & 向量寄存器位宽（128位）、元素位宽（8位）、一次运算的元素个数（8）\\
DUT、EMU & 被测设计，此处即香山处理器；EMU是由其RTL生成的仿真程序\\
NEMU & 香山配套的指令集参考模型，按指令语义计算正确结果\\
DiffTest & 每条指令提交时比较EMU与NEMU架构状态的协同仿真框架\\
AM & 运行测试程序的裸机环境\\
FU、父级 & 功能单元；“父级”指VDot所在的向量执行单元VFEX4\\
P槽、R槽 & 两槽实现中保存乘积的槽和保存结果的槽\\
ROB & 重排序缓冲区，处理器借助它按程序顺序提交指令\\
时延、发起间隔 & 一笔请求从被接收到给出结果的拍数；相邻两笔请求被接收的间隔\\
C8 & 在8行GEMV计算之间复用向量配置的函数版本\\
\bottomrule
\end{tabular}
\end{table}
